% !TeX program = pdflatex
% =============================================================================
% Arbeitsblatt 4: Schaltungen
% UZH Praktikum I -- Sara Romer Urban, Kantonsschule im Lee, HS2026
% Klasse: 1f (23.09.2026), aus PLAN_Dokumentengenerierung.md.
%
% ENTWURF. Anleitung zur praktischen Multimeter-Messung (Spannung, Strom,
% Widerstand), erstellt auf Wunsch von Loris anhand der Fotos in
% Bilder/Multimeter_Messung/ (Multimeter_Spannung.jpg, Multimeter_Strom.jpg,
% Multimeter_Widerstand.jpg) und inhaltlich abgeglichen mit "Multimeter --
% Elektrotechnik einfach erklaert.pdf" im selben Ordner (Vorbereitung/
% Messvorgang-Struktur pro Messart, Sicherheitshinweise). Jede der drei
% Messarten hat einen eigenen, selbst gezeichneten Schaltplan (circuitikz)
% mit dem Messgeraet als Kreis mit V/A/Omega, an genau der Stelle, an der es
% eingebaut wird -- wie von Loris gewuenscht, statt nur der Fotos.
%
% Knuepft an u-i-r/arbeitsblatt3-u-i-r.tex an (Einheiten, Ohmsches Gesetz)
% und an das bereits vorhandene Vorwissen laut Aufgabenstellung (Wasser-
% kreislaufanalogie, Grundlagen Multimeter-Messung: Strom in Reihe, Spannung
% parallel, DC vs. AC) -- deshalb hier keine erneute Herleitung dieser
% Grundlagen, nur die praktische Anwendung mit einem echten Geraet.
%
% Praktischer Versuch am Schluss nutzt das tatsaechlich vorhandene
% Gruppenmaterial (Netzteil 0-12V in 0.5V-Schritten, Widerstaende 47/100/
% 200 Ohm, 3 Multimeter pro Gruppe -- ermoeglicht gleichzeitige Spannungs-
% und Strommessung ohne Umstecken).
%
% -----------------------------------------------------------------------------
% Korrektur (22.09.2026, auf Wunsch): Vier Anpassungen an den Schaltplan-
% Skizzen von Spannungs-/Strommessung und am Text, nachdem Loris die
% Original-Fotos genauer eingeordnet hat: (1) Das Bauteil in den Fotos
% Multimeter_Spannung.jpg/Multimeter_Strom.jpg ist keine Widerstand,
% sondern eine Gluehbirne -- beide Schaltplaene zeigen jetzt "to[lamp]"
% statt "to[R,l=$R$]", Fliesstext/Bildunterschriften entsprechend auf
% "Gluehbirne" umgestellt (Widerstandsmessung-Abschnitt bleibt unveraendert,
% dort zeigt das Foto tatsaechlich einen Widerstand). (2) Pluspol oben,
% Minuspol unten am Netzteil -- beide Schaltplaene zeigen das jetzt explizit
% als "+"/"-"-Beschriftung neben der Batteriesymbol. (3) "Potenzial" ist den
% SuS noch nicht bekannt (nur Spannung/Strom/Widerstand) -- aus der
% Polaritaets-Hinweisbox entfernt, ersetzt durch direkten Bezug auf
% Pluspol/Minuspol wie in der Skizze gezeigt. (4) Beim eingesetzten
% Fluke-Multimeter muss nach dem Drehschalter zusaetzlich der gelbe Knopf
% oben links gedrueckt werden, um von Wechsel- auf Gleichstrom zu wechseln
% -- als Schritt in "Strom messen" ergaenzt.
%
% -----------------------------------------------------------------------------
% Erweiterung (22.09.2026, auf Wunsch): (1) Neue Einleitungsbox "Die
% Bauelemente auf Ihrem Tisch" direkt nach den Lernzielen, mit den drei
% Fotos Spannungsquelle.jpg, Steckteile_vonoben.jpg und Multimeter.jpg
% (neu in Bilder/Multimeter_Messung/) -- stellt PowerModule/Steckerteile/
% Multimeter als Bauelemente des Stromkreises vor, bevor irgendetwas
% gemessen wird. Darin neu eine rote "warnbox" (eigene Farbe, staerker als
% die orange Hinweisbox) zum Kurzschluss: Spannungsquelle nie ohne
% Widerstand/Verbraucher direkt verbinden. (2) Die Anleitungen in
% "Spannung messen"/"Strom messen"/"Widerstand messen" sind jetzt
% Bulletpoints (itemize) statt Fliesstext -- nur die Hinweisboxen bleiben
% als Fliesstext, da sie Erklaerungen sind, keine Schritt-Anleitung. (3)
% "+"/"-" an der Spannungsquelle jetzt direkt neben dem Batteriesymbol
% selbst platziert (vorher an den Ecken der Anschlussdraehte, zu weit weg
% vom eigentlichen Symbol). (4) Neue kurze Einzelaufgabe "Widerstand der
% Gluehbirne berechnen" nach "Strom messen": Ohmsches Gesetz $R=U/I$
% gegeben, Loesung mit Beispielwerten $U=4\,$V, $I=0.4\,$A $\to
% R=10\,\Omega$. (5) Der bisherige "Praktischer Versuch"-Abschnitt (ein
% einzelner Widerstand, U/I/R einzeln) ersetzt durch "Anwendung:
% Serieschaltung und Parallelschaltung": SuS bauen die drei vorhandenen
% Widerstaende (47/100/200 Ohm) einmal in Serie, einmal parallel auf und
% messen U und I an jedem Widerstand sowie $U_0$ (und bei Parallel
% zusaetzlich $I_0$), mit selbst gezeichneten circuitikz-Schaltplaenen fuer
% die Aufgabe und den echten Fotos Serieschaltung.jpg/Parallelschaltung.jpg
% (neu in Bilder/) plus den tatsaechlichen Messwerten von Loris (aus
% Messungen_Serie_Parallelschaltung.xlsx im selben Ordner) in der
% Loesungsversion. Abschliessende Vergleichsfrage: Summe der Einzel-
% spannungen (Serie) bzw. Einzelstroeme (Parallel) stimmt mit $U_0$ bzw.
% $I_0$ ueberein.
%
% -----------------------------------------------------------------------------
% Erweiterung (22.09.2026, auf Wunsch): Die bisherige Vergleichsfrage am
% Schluss ist jetzt "Kirchhoffsche Regeln herleiten" -- zusaetzlich zum
% numerischen Vergleich (Summe = Gesamtgroesse) verallgemeinern die SuS auf
% $n$ Widerstaende (bewusst nur diese einfache Verallgemeinerung ueber die
% Anzahl, keine formale Ladungserhaltungs-/Energieerhaltungs-Herleitung),
% und die Loesung benennt die beiden Regeln explizit (Maschenregel/
% Knotenregel). Danach drei neue "Merksatz"-formelboxen (Serieschaltung,
% Parallelschaltung, Schaltungen allgemein) in einer eigenen theoriebox nach
% der Gruppenarbeit: Formulierungen direkt aus Loris' Vorgabe (Wasser/Strom
% an einer Verzweigung: genau so viel kommt hinzu wie wegfliesst; die
% angelegte Spannung wird von den "Wandlern" umgewandelt; Wasser/Strom auf
% gleicher Hoehe/Spannung hat dasselbe "Wirkungspotential", daher gleiche
% Spannung an parallelen Zweigen).
%
% -----------------------------------------------------------------------------
% Korrektur/Erweiterung (22.09.2026, auf Wunsch): (1) Die "battery1"-
% Symbole waren falsch gepolt: Per Default liegt der laengere Strich
% (Pluspol) am Start des circuitikz-"to"-Pfads, bei uns also unten -- mit
% der "invert"-Option jetzt korrekt oben (deckt sich mit den textuellen
% "+"/"-"-Beschriftungen von vorhin, betrifft alle vier Schaltplaene). (2)
% Die Gluehbirne-Widerstands-Aufgabe gibt jetzt das Ohmsche Gesetz in der
% Form $U=R\cdot I$ vor (wie in u-i-r/arbeitsblatt3-u-i-r.tex), SuS stellen
% selbst nach $R$ um (Schritt gezeigt in der Loesung), statt dass $R=U/I$
% direkt vorgegeben wird. (3) Die Messwerttabellen (Serie/Parallel) sind
% jetzt geraeumiger formatiert: groesserer Zeilenabstand
% (\arraystretch{2.2}), Spaltenkoepfe mit Einheit ("$U$ in V"), breitere
% Schreiblinien (3cm), Gesamtspannung/-stromstaerke in eigener Zeile mit
% mehr Weissraum statt gedraengt neben der Tabelle. (4) Die beiden
% Schaltplaene der Anwendungsaufgabe zeigen jetzt selbst einen Kreis mit
% $V$ (parallel zu den Widerstaenden) und $A$ (in Reihe), analog zur
% Konvention aus dem Multimeter-Teil -- vorher zeigten sie nur die
% Bauteile ohne Messpunkte. (5) Serie und Parallel haben jetzt je eine
% eigene "Konklusion"-Frage direkt nach ihrer Messtabelle/Loesung (statt
% nur der gemeinsamen Schlussfrage): Konklusion Serie fragt nach dem
% Vergleich der drei Stromwerte und der Spannungssumme, Konklusion
% Parallel nach dem Vergleich der drei Spannungswerte und der Stromsumme.
% Die bisherige "Kirchhoffsche Regeln herleiten"-Frage heisst jetzt
% "Kirchhoffsche Regeln verallgemeinern" und wiederholt die numerischen
% Vergleiche nicht mehr (stehen ja schon in den beiden Konklusionen),
% sondern verallgemeinert direkt auf $n$ Widerstaende und benennt die
% beiden Regeln. (6) Der Serieschaltung-Aufgabentext verriet bisher selbst
% schon "bei einer Serieschaltung [ist der Strom] ueberall gleich gross"
% -- entfernt, SuS messen jetzt einfach an drei Stellen und sollen die
% Erkenntnis erst in der Konklusion-Frage selbst ziehen.
%
% -----------------------------------------------------------------------------
% Korrektur/Erweiterung (22.09.2026, auf Wunsch): (1) Die "Kirchhoffsche
% Regeln verallgemeinern"-Frage (oben, Punkt (6) der letzten Notiz) wieder
% zurueckgebaut: Keine Verallgemeinerung auf $n$ Widerstaende mehr (SuS
% verstehen das noch nicht), stattdessen "Kirchhoffsche Regeln benennen"
% -- bleibt konkret bei den drei gemessenen Widerstaenden
% ($U_0=U_1+U_2+U_3$ bzw. $I_0=I_1+I_2+I_3$) und fragt nur noch nach den
% Namen der beiden Regeln. Die drei Merksatz-Boxen wurden entsprechend
% von "...+U_n"/"...+I_n" auf die konkrete Dreier-Summe zurueckgestellt.
% (2) Die $V$/$A$-Messpunkt-Kreise in den beiden Schaltplaenen der
% "Anwendung"-Aufgabe (oben, Punkt (4) der letzten Notiz) wieder entfernt
% -- Schaltplaene zeigen jetzt wieder nur Batterie/Schalter/Widerstaende,
% wie vor jener Aenderung. (3) Neuer grosser Abschnitt "Ersatzwiderstand"
% nach den Merksaetzen: kurze Theoriebox zum Begriff, dann eine
% Gruppenaufgabe "Serie und Parallel" (beide Male OHNE Netzteil, direkt
% mit dem Ohmmeter gemessen -- eigene, battery-freie Schaltskizzen mit
% $\Omega$-Kreis -- Frage nach Trend bei mehr Widerstaenden sowie nach der
% Formel, SuS leiten $R_{ers}=R_1+R_2+R_3$ bzw.
% $1/R_{ers}=1/R_1+1/R_2+1/R_3$ selbst her), dann "Mischschaltung"
% ($R_1$ in Serie mit $R_2\|R_3$, ebenfalls ohne Netzteil, berechnen und
% zur Kontrolle auf dem Steckbrett messen), und zum Schluss
% "Schaltungs-Tausch": Jede Gruppe zeichnet selbst eine komplexere
% Schaltung mit mind. vier frei gewaehlten Widerstaenden, reicht sie an
% eine andere Gruppe weiter, die den Gesamtwiderstand berechnet -- bewusst
% ohne Loesung in diesem Dokument, da die Schaltung/Werte von den SuS
% selbst stammen.
%
% -----------------------------------------------------------------------------
% Korrektur (22.09.2026, auf Wunsch): In allen fuenf Schaltskizzen (beide
% "Anwendung"-Schaltplaene und alle drei "Ersatzwiderstand"-Schaltplaene)
% zeigen die Widerstaende jetzt nur noch $R_1$/$R_2$/$R_3$ als
% Bauteilbeschriftung, nicht mehr den Ohm-Wert direkt im Schaltplan. Die
% Zuordnung $R_1=\SI{47}{\ohm}$, $R_2=\SI{100}{\ohm}$, $R_3=\SI{200}{\ohm}$
% wird dafuer einmal, zentral, in der Einleitung der "Anwendung"-
% Gruppenaufgabe festgelegt (dort, wo die drei Widerstaende zum ersten Mal
% gemeinsam eingefuehrt werden) und gilt ab dort fuer den Rest des
% Dokuments. Passend dazu zeigen auch die vier Messwerttabellen
% (Anwendung Serie/Parallel, je Aufgaben- und Loesungsversion) jetzt
% $R_1$/$R_2$/$R_3$ als Zeilenbeschriftung statt des Ohm-Werts. Die
% Wiederholung der Werte in der Ersatzwiderstand-Serieschaltung-Aufgabe
% (vorher redundant zur neuen zentralen Definition) wurde entsprechend
% gekuerzt. Unveraendert bleiben: die "Steckerteile"-Bildunterschrift in
% "Die Bauelemente auf Ihrem Tisch" (beschreibt die real aufgedruckten
% Ohm-Werte auf dem Foto, keine Schaltskizze) und die generische
% "$R$"-Beschriftung im "Widerstand messen"-Abschnitt (dort noch kein
% Bezug zu den drei spezifischen Widerstaenden).
%
% -----------------------------------------------------------------------------
% Korrektur/Erweiterung (22.09.2026, auf Wunsch): (1) Die Formeln
% $U_0=U_1+U_2+U_3$ bzw. $I_0=I_1+I_2+I_3$ stehen jetzt direkt in den
% Loesungen von "Konklusion: Serieschaltung" bzw. "Konklusion:
% Parallelschaltung" (bei der jeweiligen Schaltung selbst), nicht mehr
% erst in einer eigenen Frage danach. (2) Die bisherige Frage
% "Kirchhoffsche Regeln benennen" (vorher Aufgabe 6) ist deshalb komplett
% entfernt -- die Namen der beiden Regeln (Knotenregel/Maschenregel)
% stehen weiterhin unveraendert im Merksatz "Schaltungen allgemein".
% Alle \setcounter{question}-Werte der nachfolgenden Gruppenaufgaben
% (Ersatzwiderstand Serie/Parallel, Mischschaltung, Schaltungs-Tausch) um
% eins nach unten korrigiert. (3) Die Formel-Anzeigen in den Merksatz-
% Boxen "Serieschaltung" und "Parallelschaltung" entfernt (nur noch
% Fliesstext ohne die $U_0=\dots$/$I_0=\dots$-Gleichung) -- die Formeln
% stehen ja jetzt schon bei den Konklusionen (Punkt (1)). (4) Die Loesung
% von "Ersatzwiderstand: Parallelschaltung" zeigt jetzt den genauen
% Rechenweg (Bruch mit den drei Widerstaenden aufschreiben, jeden
% Summanden einzeln als Dezimalzahl ausrechnen, addieren, dann den
% Kehrwert bilden) statt direkt von der Summenformel zum Endergebnis zu
% springen. (5) "Schaltungs-Tausch" hat jetzt ein Beispiel mit
% vollstaendiger Loesung ($R_a=100\,\Omega$, $R_b=200\,\Omega$ in Serie,
% parallel zu $R_c=300\,\Omega$, das Ganze in Serie mit $R_d=50\,\Omega$,
% Ergebnis $\SI{200}{\ohm}$) -- eigene Schaltskizze, offen sichtbar als
% Anschauungsbeispiel, die Rechnung selbst in einer \solution-Box.
% Korrektur (22.09.2026, auf Wunsch): Das Beispiel selbst (Schaltskizze
% plus Kontext-Satz) steckt jetzt ebenfalls in der \solution-Box, nicht
% mehr offen sichtbar davor -- Loris meinte mit "Beispiel mit Loesung"
% ein Beispiel, das komplett zur Loesung gehoert (also nur im
% Lehrpersonen-/\printanswers-Exemplar sichtbar), nicht ein zusaetzliches
% offenes Anschauungsbeispiel vor der eigentlichen Aufgabe.
% Ausserdem (auf Wunsch): \Thema und \ArbeitsblattTitel von "Messen mit
% dem Multimeter" auf "Schaltungen" umbenannt (Dokument deckt inzwischen
% deutlich mehr als nur die Multimeter-Messung ab: Ersatzwiderstand,
% Schaltungs-Tausch, UND-/ODER-Schaltung) -- betrifft Kopfzeile,
% Seitenkopf und PDF-Metadaten (\pdftitle), Dateiname bleibt unveraendert
% (arbeitsblatt4-schaltungen.tex passt bereits).
%
% -----------------------------------------------------------------------------
% Erweiterung (22.09.2026, auf Wunsch): Neue Zusatzaufgabe fuer schnelle
% Gruppen ganz am Schluss: "UND- und ODER-Schaltung", zwei Schalter plus
% eine Gluehbirne. Anders als die bisherigen Aufgaben bekommen die SuS
% hier den Schaltplan NICHT vorgegeben -- sie sollen ihn nach dem Bauen
% selbst zeichnen (\Antwortraum statt fertigem circuitikz-Schaltplan in
% der Aufgabe selbst). Die Loesung zeigt je Loris' eigenes, tatsaechlich
% gebautes Referenzfoto (Bilder/schaltung_und.jpg: zwei Schalter in Serie
% vor der Gluehbirne -- Gluehbirne leuchtet nur, wenn beide Schalter zu
% sind; Bilder/schaltung_oder.jpg: zwei Schalter parallel -- ein
% einzelner Schalter genuegt) sowie einen eigenen circuitikz-
% Referenzschaltplan zum Abgleich.
%
% -----------------------------------------------------------------------------
% Review-Korrektur (22.09.2026, auf Wunsch -- Redundanzen/Logik/
% physikalische Korrektheit von Arbeitsblatt+Aufgabenblatt gemeinsam
% geprueft, physics-accuracy-reviewer und redundanz-tiefe-reviewer als
% Subagenten sowie eigene Durchsicht):
% (1) [physics-accuracy, behoben] Ersatzwiderstand-Mischschaltung-Loesung
% sprang beim Berechnen von $R_{23}=R_2\|R_3$ direkt von der Bruchsumme
% zum Endergebnis (ein einziger \Longrightarrow-Schritt fuer
% Gleichnenner-Bildung UND Kehrwert zugleich) -- jetzt in drei einzelne
% Schritte aufgeteilt (Bruchsumme aufschreiben, gemeinsamer Nenner, dann
% erst Kehrwert), konsistent mit jeder anderen Bruchrechnung in beiden
% Dokumenten. (2) [physics-accuracy, BLOCKING, behoben] Aufgabenblatt-
% Aufgabe 3a verlangte als Begruendung implizit die allgemeine Regel
% "Ersatzwiderstand einer Parallelschaltung ist immer kleiner als der
% kleinste Einzelwiderstand" -- diese Regel stand bisher nirgends
% sichtbar fuer SuS, nur einmal verdeckt in einer \solution-Box fuer
% einen einzelnen Zahlenfall. Neue sichtbare Box "Merksatz:
% Ersatzwiderstand" im Arbeitsblatt (nach der Ersatzwiderstand-
% Gruppenaufgabe, vor der Mischschaltung) haelt jetzt explizit fest:
% Serie-Ersatzwiderstand immer groesser als jeder Einzelwiderstand,
% Parallel-Ersatzwiderstand immer kleiner als der kleinste
% Einzelwiderstand -- inkl. der beiden Formeln. (3) [eigener Fund]
% Die Ersatzwiderstand-Serie/Parallel-Loesungen behaupteten faelschlich
% "Gemessen: $R_{ers}\approx...$", obwohl es sich um den theoretisch
% berechneten Wert handelt (keine echten Messdaten wie bei der
% Anwendung-Aufgabe) -- umformuliert zu "Ihr gemessener Wert sollte damit
% uebereinstimmen", konsistent mit der bereits korrekten Formulierung bei
% der Mischschaltung. (4) [redundanz-tiefe, behoben] Die dritte
% Merksatz-Box "Schaltungen allgemein" wiederholte die Kernsaetze aus den
% Boxen "Serieschaltung"/"Parallelschaltung" fast woertlich und fuegte
% nur die Regelnamen hinzu -- auf einen kurzen Rueckverweis gekuerzt, der
% nur noch die Namen Knotenregel/Maschenregel an die bereits
% festgehaltenen Regeln haengt, ohne sie erneut zu definieren.
% Im Aufgabenblatt: [physics-accuracy, MINOR, behoben] Aufgabe 3c
% (Teilspannungen) verwendete den auf 0.102 A gerundeten Strom aus 3b
% weiter, was zu einem 1-Rappen-Rundungsfehler bei $U_3$ fuehrte
% (8.16 V statt der praeziseren 8.17 V) -- Zwischenwert in 3b auf
% 0.1021 A praezisiert, 3c entsprechend auf 8.17 V/3.83 V angepasst.
% Alle anderen von beiden Reviewern geprueften Punkte (Zahlenwerte,
% Schaltplan-Topologien, Aufgabenschwierigkeit/-progression,
% Diagrammpaare) waren bereits korrekt bzw. absichtlich so gestaltet --
% keine weiteren Aenderungen noetig.
%
% -----------------------------------------------------------------------------
% Umbau (22.09.2026, auf Wunsch): (1) Der Ersatzwiderstand wird jetzt
% direkt bei "Anwendung: Serieschaltung und Parallelschaltung" mitgemessen,
% statt in einer eigenen, spaeteren Gruppenaufgabe dieselbe Schaltung noch
% einmal aufzubauen. Die SuS haben den Stromkreis ja schon vor sich.
% Jede der beiden Teilaufgaben (Serie/Parallel) hat dafuer einen neuen
% Schritt: Spannungsquelle trennen (Sicherheitsregel fuer Widerstands-
% messungen, siehe "Widerstand messen"), dann $R_{ges}$ direkt mit dem
% Ohmmeter messen; die Messwerttabelle hat dafuer eine zusaetzliche Zeile.
% Die beiden "Konklusion"-Fragen fragen jetzt zusaetzlich nach dem Trend
% (groesser/kleiner) und nach der Formel fuer $R_{ers}$, analog zur
% bisherigen Frage in der separaten Gruppenaufgabe. Die eigenstaendige
% Gruppenaufgabe "Ersatzwiderstand: Serie und Parallel" (mit eigenem
% Schaltungsaufbau ohne Netzteil) sowie die vorangehende kurze
% "Ersatzwiderstand"-Theoriebox entfallen dadurch komplett; die
% Definition des Ersatzwiderstands steckt jetzt in einem Satz in der Einleitung
% der "Anwendung"-Gruppenaufgabe. (2) Die Box "Merksatz: Ersatzwiderstand"
% (zwei Formelboxen mit den fertigen Formeln $R_{ers}=R_1+R_2+R_3$ bzw.
% $1/R_{ers}=1/R_1+1/R_2+1/R_3$) ist ebenfalls entfernt: Das Herleiten
% dieser Formeln ist Teil der Aufgabe (Konklusion-Fragen oben), sie sollen
% den SuS nicht vorher schon vorgegeben werden. Die "Ersatzwiderstand:
% Mischschaltung"-Aufgabe verweist weiterhin auf "die Formeln von oben":
% Gemeint sind jetzt die Formeln, die die SuS selbst in ihren eigenen
% Konklusion-Antworten hergeleitet haben. (3) Alle nachfolgenden
% \setcounter{question}-Werte (Mischschaltung, Schaltungs-Tausch,
% Zusatzaufgabe) um zwei nach unten korrigiert, da die entfallene
% Gruppenaufgabe zwei Fragennummern belegt hatte. (4) Die Lernzielbox
% deckt jetzt knapp das gesamte Dokument ab (Multimeter-Messung inkl.
% Sicherheit, Serie-/Parallelverhalten, Ersatzwiderstand) statt nur die
% Multimeter-Messung. Die alte dritte Lernzielzeile zur Sicherheit ist
% dafuer in die erste eingeflossen, um bei drei Punkten zu bleiben.
%
% -----------------------------------------------------------------------------
% Korrektur (23.09.2026, auf Wunsch): (1) "Strom messen": Der gelbe Knopf
% zum Umschalten auf Gleichstrom sitzt am tatsaechlichen Fluke-Multimeter
% oben rechts, nicht oben links -- Text korrigiert. (2) Die
% Kurzschluss-Warnbox erklaert jetzt auch WARUM ein Kurzschluss gefaehrlich
% ist, nicht nur DASS er es ist: Nach dem Ohmschen Gesetz $I=U/R$ wird die
% Stromstaerke riesig, wenn der Widerstand fast null ist. (3) "Widerstand
% messen": Der Hinweis zum Trennen des Widerstands vom Stromkreis erwaehnte
% bisher "auslöten bzw. ausstecken" -- am Steckbrett wird aber nichts
% geloetet, nur umgesteckt, deshalb "auslöten" gestrichen (im ganzen
% Dokument geprueft, war die einzige Stelle). (4) Mehrere Aufgaben stellten
% bisher zwei oder drei inhaltlich verschiedene Fragen hintereinander vor
% einer einzigen gemeinsamen Antwortbox: "Konklusion: Serieschaltung" und
% "Konklusion: Parallelschaltung" (je drei Fragen: Stromvergleich,
% Spannungssumme, Ersatzwiderstand-Formel) sowie "Schaltungs-Tausch" (Skizze
% zeichnen und Widerstand berechnen waren zwei verschiedene Antwortraeume,
% aber ohne \parts-Struktur). Alle drei jetzt in \parts mit je einer Frage
% und einer eigenen Antwortbox pro \part umgebaut, damit immer nur eine
% Frage auf einmal gestellt wird. Die Loesungs-Inhalte entsprechend auf die
% einzelnen \part-Antworten aufgeteilt (bei Schaltungs-Tausch bleibt die
% Gesamtloesung als ein zusammenhaengendes Beispiel nach den \parts stehen,
% da sie eine einzige durchgehende Beispielrechnung zeigt).
%
% -----------------------------------------------------------------------------
% Korrektur (23.09.2026, auf Wunsch): (1) Bei "Konklusion: Serieschaltung"
% und "Konklusion: Parallelschaltung" gibt der letzte \part jetzt die
% Formel fuer den Ersatzwiderstand direkt vor ($R_{ers}=R_1+R_2+R_3$ bzw.
% $1/R_{ers}=1/R_1+1/R_2+1/R_3$), statt sie selbst herleiten zu lassen --
% die SuS ueberpruefen die Formel stattdessen mit ihrem gemessenen $R_{ges}$
% und den drei Einzelwiderstaenden. Die Loesungen entsprechend umformuliert
% (erst einsetzen/ausrechnen, dann mit dem Messwert vergleichen, statt
% umgekehrt). (2) Die Lernziele-Box am Dokumentanfang komplett entfernt.
%
% -----------------------------------------------------------------------------
% Korrektur (23.09.2026, auf Wunsch): Bei "Konklusion: Serieschaltung" (c)
% wieder zurueckgebaut auf die vorherige Fassung: Die SuS stellen die
% Formel fuer $R_{ers}$ wieder selbst auf, statt sie vorgegeben zu
% bekommen und nur zu ueberpruefen (Punkt (1) der letzten Notiz). Bei
% "Konklusion: Parallelschaltung" (c) bleibt die Formel dagegen vorgegeben
% -- dort ist das Herleiten (Kehrwertsumme) deutlich anspruchsvoller als
% die einfache Addition bei der Serieschaltung.
% =============================================================================
\documentclass[addpoints,11pt,a4paper]{exam}

\usepackage[T1]{fontenc}
\usepackage[utf8]{inputenc}
\usepackage[ngerman]{babel}
\usepackage{lmodern}
\usepackage[a4paper,margin=2.2cm,headheight=14pt]{geometry}
\usepackage{amsmath}
\usepackage{siunitx}
% Hausstil: zusammengesetzte Einheiten immer als echter Bruch (\frac), nicht
% als Inline-Schraegstrich oder \tfrac -- gilt global fuer jedes \si/\SI.
\sisetup{per-mode=fraction,fraction-command=\frac}
\usepackage{array,booktabs,tabularx,multirow}
\usepackage{enumitem}
\usepackage{xcolor}
\usepackage{colortbl}
\usepackage{tikz}
\usepackage[european]{circuitikz}
\usepackage{physics}
\usepackage[most]{tcolorbox}
\usepackage{lastpage}
\usepackage{graphicx}
\usepackage{hyperref}

\usetikzlibrary{arrows.meta,calc,angles,quotes}

\usepackage{setspace}
\usepackage{needspace}
\setlength{\parindent}{0pt}
\setlength{\parskip}{0.6em}
\setstretch{1.08}
\setlist[itemize]{itemsep=0.4em}

% -----------------------------
% Farben (Corporate-Farbschema Physik KFR)
% -----------------------------
\definecolor{titleblue}{HTML}{183A6B}
\definecolor{accentblue}{HTML}{2E6F9E}
\definecolor{lightblue}{HTML}{EAF4FB}
\definecolor{lightgray}{HTML}{F4F6F8}
\definecolor{darkgray}{HTML}{4A4A4A}
\definecolor{accentorange}{HTML}{D9720C}
\definecolor{accentpurple}{HTML}{7B2D8E}
\definecolor{lightorange}{HTML}{FCEEE0}
\definecolor{lightpurple}{HTML}{F3EAF7}

% Links (hyperref): farblich ans Corporate-Farbschema angeglichen.
\hypersetup{
  colorlinks=true,
  urlcolor=accentblue,
  linkcolor=accentblue,
  pdftitle={Arbeitsblatt 4: Schaltungen},
  pdfauthor={Loris Keller}
}

% -----------------------------
% Boxen
% -----------------------------
\tcbset{before skip=10pt, after skip=10pt}
\newtcolorbox{titlebox}{
  enhanced,
  colback=titleblue,
  colframe=titleblue,
  boxrule=0pt,
  arc=3mm,
  left=3mm,right=3mm,top=3mm,bottom=3mm
}

\newlength{\aufgabenindent}
\settowidth{\aufgabenindent}{10.\hskip\labelsep}

% Praktikum I: Lernziele/Einleitung/Theorie/Aufgaben werden NICHT farbig
% gerahmt (nur Formel-, Hinweis- und Antwortbox bleiben Boxen) -- siehe
% institutionen/UZH-Praktikum-I/CLAUDE.md, Abschnitt "Boxen-Layout".
\newenvironment{infobox}{%
  \par\addvspace{10pt}\noindent
}{%
  \par\addvspace{10pt}
}

% Titel von Theorie-/Aufgaben-/Gruppenarbeit-Abschnitten: groesser, farbig
% und mit duenner Trennlinie statt schlichtem Fettdruck.
\newcommand{\Teiltitel}[2]{%
  \noindent{\Large\bfseries\color{#1}#2}\par
  \vspace{1pt}\noindent{\color{#1}\rule{\linewidth}{0.8pt}}\par\smallskip\noindent
}

\newenvironment{theoriebox}[1][Theorie]{%
  \par\addvspace{12pt}\Teiltitel{titleblue}{#1}
}{%
  \par\addvspace{10pt}
}

\newtcolorbox{taskbox}[1][]{
  enhanced, breakable,
  colback=white, colframe=white, boxrule=0pt,
  left=0mm,right=0mm,top=0mm,bottom=0mm,
  before skip=12pt, after skip=10pt,
  before upper={%
    \def\tmpboxtitle{#1}%
    \ifx\tmpboxtitle\empty\else\Teiltitel{accentblue}{#1}\fi
  }
}

\newtcolorbox{groupbox}[1][Gruppenarbeit]{
  enhanced, breakable,
  colback=white, colframe=white, boxrule=0pt,
  left=0mm,right=0mm,top=0mm,bottom=0mm,
  before skip=12pt, after skip=10pt,
  before upper={\Teiltitel{accentorange}{#1}}
}

\newtcolorbox{hinweisbox}[1][Hinweis]{
  enhanced,
  breakable,
  colback=lightorange,
  colframe=accentorange,
  title={#1},
  fonttitle=\bfseries,
  boxrule=0.7pt,
  arc=2mm,
  left=5mm,right=5mm,top=2mm,bottom=2mm
}

% Warnbox: eigene rote Farbe, deutlich staerker als die orange Hinweisbox --
% fuer Sicherheitswarnungen mit Beschaedigungs-/Verletzungsrisiko (hier:
% Kurzschluss), nicht fuer gewoehnliche Verfahrenshinweise.
\definecolor{accentred}{HTML}{B3261E}
\definecolor{lightred}{HTML}{FBE9E7}
\newtcolorbox{warnbox}[1][Warnung]{
  enhanced,
  breakable,
  colback=lightred,
  colframe=accentred,
  title={#1},
  fonttitle=\bfseries,
  boxrule=0.9pt,
  arc=2mm,
  left=5mm,right=5mm,top=2mm,bottom=2mm
}

\newtcolorbox{formelbox}[1][Formel]{
  enhanced,
  breakable,
  colback=lightpurple,
  colframe=accentpurple,
  title={#1},
  fonttitle=\bfseries,
  boxrule=0.9pt,
  arc=2mm,
  left=5mm,right=5mm,top=2mm,bottom=2mm
}

\newtcolorbox{answerbox}{
  breakable,
  colback=white,
  colframe=gray!55,
  boxrule=0.5pt,
  arc=1.5mm,
  left=2mm,right=2mm,top=2mm,bottom=2mm
}

% Bild-Helfer: zentriertes Bild mit spuerbarem Abstand zum umgebenden Text.
% #1 = Breite, #2 = Dateipfad.
\newcommand{\Bild}[2]{%
  \par\vspace{0.6cm}
  \begin{center}
    \includegraphics[width=#1]{#2}
  \end{center}
  \vspace{0.6cm}
}

% --- Loesungen ein-/ausblenden -----------------------------------------------
\noprintanswers
%\printanswers

\pointformat{}
\renewcommand{\solutiontitle}{\noindent\color{titleblue}\textbf{L\"osung:}\enspace}

\pagestyle{headandfoot}
\cfoot{\textcolor{darkgray}{Seite \thepage\ von \pageref{LastPage}}}

\newcommand{\Kopfzeile}[4]{%
  \begin{titlebox}
    {\color{white}\sffamily\bfseries\Large #4}
  \end{titlebox}
  \vspace{0.5em}
  \noindent\textbf{Klasse:}~#1 \hfill \textbf{Datum:}~#2 \hfill \textbf{Thema:}~#3
    \hfill \textbf{Lehrperson:}~Loris Keller
  \vspace{0.8em}
  \lhead{\textcolor{darkgray}{#3}}
  \rhead{\textcolor{darkgray}{#4}}
}

\newlength{\antwortraumlen}
\newcommand{\Antwortraum}[1]{%
  \pgfmathsetlength{\antwortraumlen}{#1*2.5cm}%
  \begin{answerbox}
    \rule{0pt}{\antwortraumlen}%
  \end{answerbox}%
}

% Werte fuer Kopfzeile zentral setzen
\newcommand{\Klasse}{1f}
\newcommand{\Datum}{23.09.2026}
\newcommand{\Thema}{Schaltungen}
\newcommand{\ArbeitsblattTitel}{Arbeitsblatt 4: Schaltungen}

\begin{document}

\Kopfzeile{\Klasse}{\Datum}{\Thema}{\ArbeitsblattTitel}

% =============================================================================
% Bauelemente des Stromkreises
% =============================================================================
\needspace{10cm}
\begin{theoriebox}[Die Bauelemente auf Ihrem Tisch]
Bevor Sie messen, ein kurzer Überblick über die drei Geräte, mit denen Sie
heute arbeiten:

\Bild{6cm}{Bilder/Multimeter_Messung/Spannungsquelle.jpg}
{\scriptsize \textbf{Spannungsquelle} (PowerModule): liefert die Spannung
für Ihren Stromkreis, einstellbar von \SI{0}{\volt} bis \SI{12}{\volt} in
\SI{0.5}{\volt}-Schritten. Oben der Pluspol, unten der Minuspol.}

\Bild{9cm}{Bilder/Multimeter_Messung/Steckteile_vonoben.jpg}
{\scriptsize \textbf{Steckerteile}: Sie stecken diese Bauteile auf das
Steckbrett, um Ihren Stromkreis aufzubauen. Zu sehen sind die Widerstände
$\SI{47}{\ohm}$, $\SI{100}{\ohm}$, $\SI{200}{\ohm}$, ein leeres
Verbindungsstück, ein Schalter, eine Glühbirne und eine LED, jeweils mit
ihrem Schaltzeichen beschriftet.}

\Bild{5cm}{Bilder/Multimeter_Messung/Multimeter.jpg}
{\scriptsize \textbf{Multimeter}: das Messgerät, mit dem Sie Spannung,
Strom und Widerstand messen.}

\needspace{5cm}
\begin{warnbox}[Achtung: Kurzschluss]
Verbinden Sie die Spannungsquelle \textbf{\emph{nie}} direkt mit sich
selbst, also ohne einen Widerstand oder Verbraucher (z.\,B. eine
Glühbirne) im Stromkreis. Nach dem Ohmschen Gesetz $I=U/R$ wird die
Stromstärke $I$ riesig gross, wenn der Widerstand $R$ fast null ist: Es
fliesst ein sehr grosser Strom, ein \textbf{\emph{Kurzschluss}}. Die
Spannungsquelle, die Kabel oder das Steckbrett können dabei beschädigt
werden. Bauen Sie deshalb immer zuerst Ihren vollständigen Stromkreis mit
einem Widerstand oder Verbraucher auf, bevor Sie die Spannungsquelle
einschalten.
\end{warnbox}
\end{theoriebox}

% =============================================================================
% Das Messgeraet im Schaltplan
% =============================================================================
\needspace{8cm}
\begin{theoriebox}[Das Messgerät im Schaltplan]
Sie wissen bereits: Eine Spannungsmessung wird \textbf{\emph{parallel}} zum
Bauteil angeschlossen, eine Strommessung \textbf{\emph{in Reihe}} in den
Stromkreis eingebaut. Im Schaltplan zeichnet man das Multimeter dafür
vereinfacht als \textbf{\emph{Kreis mit dem Symbol der gemessenen Grösse}}
darin: $V$ für Spannung, $A$ für Strom, $\Omega$ für Widerstand. Die Lage
dieses Kreises im Schaltplan zeigt Ihnen sofort, wie das Multimeter
tatsächlich angeschlossen werden muss.
\end{theoriebox}

% =============================================================================
% Spannung messen
% =============================================================================
\needspace{6cm}
\begin{theoriebox}[Spannung messen]
\begin{itemize}[leftmargin=1.2cm,itemsep=0.3em]
  \item Schwarzes Messkabel in die COM-Buchse, rotes Messkabel in die
  Buchse mit dem $V$-Zeichen.
  \item Drehschalter auf Gleichspannung ($\overline{V}$) stellen.
  \item Beide Prüfspitzen an die zwei Punkte halten, zwischen denen Sie
  die Spannung messen möchten, hier an den beiden Enden der Glühbirne im
  Stromkreis, \textbf{\emph{parallel}} dazu (siehe Schaltplan unten).
\end{itemize}

\begin{formelbox}[Schaltplan: Spannungsmessung]
\begin{center}
\begin{circuitikz}[scale=1.3]
  \draw (0,0) to[battery1,invert,l=$U_0$] (0,3)
              to[normal open switch] (2.2,3)
              -- (4.4,3)
              to[lamp] (4.4,0)
              -- (0,0);
  \node at (0.4,1.85) {$+$};
  \node at (0.4,1.15) {$-$};
  \draw (4.4,3) -- (6,3) -- (6,1.9);
  \node[draw,circle,minimum size=0.8cm] (V) at (6,1.5) {$V$};
  \draw (6,1.1) -- (6,0) -- (4.4,0);
\end{circuitikz}
\end{center}
{\scriptsize Das Voltmeter (Kreis mit $V$) liegt \emph{parallel} zur
Glühbirne: Es misst die Spannung, die über ihr abfällt, ohne den
Stromkreis selbst zu unterbrechen. Am Netzteil oben der Pluspol, unten
der Minuspol.}
\end{formelbox}

\Bild{7cm}{Bilder/Multimeter_Messung/Multimeter_Spannung.jpg}
{\scriptsize Spannungsmessung mit dem Multimeter: Beide Prüfspitzen liegen
parallel an den beiden Enden der Glühbirne an, hier zeigt das Display
noch $\SI{0}{\volt}$ (Stromkreis noch nicht eingeschaltet).}

\needspace{5cm}
\begin{hinweisbox}[Polarität bei Gleichspannung]
Legen Sie die rote Prüfspitze an die Seite mit dem Pluspol (in der Skizze
oben) und die schwarze Prüfspitze an die Seite mit dem Minuspol (unten).
Sie erhalten dann einen positiven Wert. Vertauschen Sie die Spitzen,
zeigt das Multimeter denselben Betrag mit negativem Vorzeichen an, z.\,B.
$\SI{-6}{\volt}$.
\end{hinweisbox}
\end{theoriebox}

% =============================================================================
% Strom messen
% =============================================================================
\needspace{6cm}
\begin{theoriebox}[Strom messen]
\begin{itemize}[leftmargin=1.2cm,itemsep=0.3em]
  \item Schwarzes Messkabel weiterhin in die COM-Buchse, rotes Messkabel
  aber in die Buchse mit dem $A$-Zeichen (mA oder A, je nach erwarteter
  Stromstärke; bei Unsicherheit zuerst die grössere Buchse verwenden).
  \textbf{\emph{Achtung:}} Für eine Strommessung braucht es eine andere
  Buchse als für Spannung und Widerstand.
  \item Drehschalter auf Gleichstrom ($\overline{A}$) stellen. Danach
  zusätzlich den gelben Knopf oben rechts drücken, um von Wechselstrom auf
  Gleichstrom umzuschalten.
  \item Stromkreis spannungsfrei schalten und an einer Stelle auftrennen.
  Multimeter \textbf{\emph{in Reihe}} in diese Lücke einsetzen (siehe
  Schaltplan unten), dann den Stromkreis wieder einschalten.
\end{itemize}

\begin{formelbox}[Schaltplan: Strommessung]
\begin{center}
\begin{circuitikz}[scale=1.8]
  \draw (0,0) to[battery1,invert,l=$U_0$] (0,3)
              to[normal open switch] (2,3)
              -- (3.2,3);
  \node at (0.4,1.85) {$+$};
  \node at (0.4,1.15) {$-$};
  \node[draw,circle,minimum size=0.8cm] (A) at (3.6,3) {$A$};
  \draw (4,3) -- (5.2,3)
              to[lamp] (5.2,0)
              -- (0,0);
\end{circuitikz}
\end{center}
{\scriptsize Das Amperemeter (Kreis mit $A$) liegt \emph{in Reihe} in der
Schleife: Durch das Messgerät fliesst genau derselbe Strom $I$ wie durch
die Glühbirne. Am Netzteil oben der Pluspol, unten der Minuspol.}
\end{formelbox}

\needspace{8.5cm}
\Bild{7cm}{Bilder/Multimeter_Messung/Multimeter_Strom.jpg}
{\scriptsize Strommessung mit dem Multimeter: Beide Prüfspitzen sind in
die aufgetrennte Schleife eingesetzt, das Multimeter ist Teil des
Stromkreises.}

\begin{hinweisbox}[Sicherheit bei der Strommessung]
Trennen Sie den Stromkreis nur, wenn er spannungsfrei ist (Netzteil aus
oder Schalter offen). Schliessen Sie danach den Stromkreis wieder, um zu
messen. Bei falscher Polarität zeigt das Multimeter denselben Betrag mit
negativem Vorzeichen an, genau wie bei der Spannungsmessung.
\end{hinweisbox}
\end{theoriebox}

% =============================================================================
% Kurze Aufgabe: Widerstand der Gluehbirne berechnen
% =============================================================================
\needspace{8cm}
\begin{questions}
\begin{taskbox}[Aufgabe: Widerstand der Glühbirne berechnen]
\question[3] Sie haben oben an derselben Glühbirne bereits die Spannung
$U$ und die Stromstärke $I$ gemessen. Das Ohmsche Gesetz lautet
\[
  U = R\cdot I.
\]
Stellen Sie die Formel selbst nach $R$ um und berechnen Sie damit den
Widerstand $R$ der Glühbirne.
\begin{answerbox}
\vspace{2.5cm}
\end{answerbox}
\begin{solution}
Umstellen (beide Seiten durch $I$ teilen):
\[
  \frac{U}{I} = \frac{R\cdot I}{I} = R \qquad\Longrightarrow\qquad R=\frac{U}{I}.
\]
Zum Beispiel bei $U=\SI{4}{\volt}$ und $I=\SI{0.4}{\ampere}$:
\[
  R = \frac{U}{I} = \frac{\SI{4}{\volt}}{\SI{0.4}{\ampere}} = \SI{10}{\ohm}.
\]
(Mit Ihren eigenen Messwerten von oben erhalten Sie einen etwas anderen
Wert.)
\end{solution}
\end{taskbox}
\end{questions}

% =============================================================================
% Widerstand messen
% =============================================================================
\needspace{6cm}
\begin{theoriebox}[Widerstand messen]
\begin{itemize}[leftmargin=1.2cm,itemsep=0.3em]
  \item Schwarzes Messkabel in die COM-Buchse, rotes Messkabel in die
  Buchse mit dem $\Omega$-Zeichen.
  \item Drehschalter auf das $\Omega$-Zeichen stellen.
\end{itemize}

\begin{hinweisbox}[Wichtig: Schaltung muss spannungsfrei sein]
Für eine Widerstandsmessung darf am zu messenden Bauteil \textbf{\emph{keine
Spannungsquelle}} angeschlossen sein. Sonst misst das Multimeter einen
falschen Wert, da es selbst einen kleinen Strom durch das Bauteil
schickt. Trennen Sie den Widerstand deshalb ganz vom Stromkreis (mindestens auf
einer Seite ausstecken), bevor Sie messen.
\end{hinweisbox}

\begin{itemize}[leftmargin=1.2cm,itemsep=0.3em]
  \item Beide Prüfspitzen direkt an die zwei Anschlüsse des Widerstands
  halten (siehe Schaltplan unten).
\end{itemize}

\begin{formelbox}[Schaltplan: Widerstandsmessung]
\begin{center}
\begin{circuitikz}[scale=2.0]
  \draw (0,0) to[R,l=$R$] (0,3);
  \draw (0,3) -- (4,3) -- (4,1.9);
  \node[draw,circle,minimum size=0.8cm] (Om) at (4,1.5) {$\Omega$};
  \draw (4,1.1) -- (4,0) -- (0,0);
\end{circuitikz}
\end{center}
{\scriptsize Der Widerstand $R$ hängt hier an \emph{keinem} Netzteil: Das
Ohmmeter (Kreis mit $\Omega$) ist direkt mit seinen beiden Anschlüssen
verbunden.}
\end{formelbox}

\Bild{9cm}{Bilder/Multimeter_Messung/Multimeter_Widerstand.jpg}
{\scriptsize Widerstandsmessung mit dem Multimeter: Der Widerstand ist
vom Steckbrett getrennt, das Display zeigt seinen Wert
($\SI{46.8}{\ohm}$, gemessen an einem als $\SI{47}{\ohm}$
beschrifteten Widerstand; kleine Abweichungen dieser Grössenordnung
sind normal).}
\end{theoriebox}

% =============================================================================
% Anwendung: Serieschaltung und Parallelschaltung
% =============================================================================
\needspace{10cm}
\begin{questions}
\setcounter{question}{1}
\begin{groupbox}[Anwendung: Serieschaltung und Parallelschaltung]
Bilden Sie Vierergruppen. Jede Gruppe hat ein Steckbrett, ein Netzteil,
einen Schalter, die drei Widerstände $R_1=\SI{47}{\ohm}$,
$R_2=\SI{100}{\ohm}$ und $R_3=\SI{200}{\ohm}$ sowie drei Multimeter. Sie
bauen nun dieselben drei Widerstände einmal in Serie und einmal parallel
auf und vergleichen, wie sich Spannung und Strom dabei verteilen. Bei
beiden Schaltungen messen Sie ausserdem deren
\textbf{\emph{Ersatzwiderstand}} $R_{ges}$: den Widerstand, den ein
einziger Widerstand haben müsste, um die ganze Schaltung zu ersetzen.

\needspace{10cm}
\question[6] \textbf{Serieschaltung}
\begin{itemize}[leftmargin=1.2cm,itemsep=0.3em]
  \item Bauen Sie die drei Widerstände gemäss Schaltplan in Serie mit
  Schalter und Netzteil auf. Stellen Sie das Netzteil auf $\SI{4}{\volt}$
  ein, schliessen Sie den Schalter aber noch nicht.
  \item Messen Sie nacheinander die Spannung über jedem der drei
  Widerstände sowie die Gesamtspannung $U_0$ über der ganzen Schaltung.
  \item Messen Sie die Stromstärke $I$ jeweils in der Nähe jedes der drei
  Widerstände (siehe Schaltplan unten) und tragen Sie alle drei Werte in
  die Tabelle ein.
  \item Trennen Sie anschliessend die Spannungsquelle vom Stromkreis:
  Für eine Widerstandsmessung darf keine Spannungsquelle angeschlossen
  sein (siehe \glqq Widerstand messen\grqq{} weiter oben). Messen Sie
  dann mit dem Multimeter ($\Omega$-Modus) direkt den Gesamtwiderstand
  $R_{ges}$ derselben Serieschaltung.
\end{itemize}
\begin{center}
\begin{circuitikz}[scale=1.15]
  \draw (0,0) to[battery1,invert,l=$U_0$] (0,3)
              to[normal open switch] (2,3)
              to[R,l=$R_1$] (4.4,3)
              to[R,l=$R_2$] (6.4,3)
              to[R,l=$R_3$] (8.4,3)
              -- (8.4,0)
              -- (0,0);
\end{circuitikz}
\end{center}
\begin{center}
\renewcommand{\arraystretch}{2.2}
\begin{tabular}{lcc}
\toprule
$R$ & $U$ in $\si{\volt}$ & $I$ in $\si{\ampere}$ \\
\midrule
$R_1$ & \rule{3cm}{0.4pt} & \rule{3cm}{0.4pt} \\
$R_2$ & \rule{3cm}{0.4pt} & \rule{3cm}{0.4pt} \\
$R_3$ & \rule{3cm}{0.4pt} & \rule{3cm}{0.4pt} \\
\bottomrule
\end{tabular}
\renewcommand{\arraystretch}{1}

\vspace{0.5cm}
Gesamtspannung $U_0 = $ \rule{3cm}{0.4pt}

\vspace{0.3cm}
Gesamtwiderstand $R_{ges} = $ \rule{3cm}{0.4pt}
\end{center}
\begin{solution}
\Bild{7cm}{Bilder/Serieschaltung.jpg}
\begin{center}
\renewcommand{\arraystretch}{2.2}
\begin{tabular}{lcc}
\toprule
$R$ & $U$ in $\si{\volt}$ & $I$ in $\si{\ampere}$ \\
\midrule
$R_1$ & $0.552$ & $0.012$ \\
$R_2$ & $1.173$ & $0.012$ \\
$R_3$ & $2.321$ & $0.012$ \\
\bottomrule
\end{tabular}
\renewcommand{\arraystretch}{1}

\vspace{0.5cm}
Gesamtspannung $U_0 = \SI{4}{\volt}$

\vspace{0.3cm}
Gesamtwiderstand $R_{ges} = \SI{347}{\ohm}$
\end{center}
\end{solution}

\needspace{6cm}
\question[5] \textbf{Konklusion: Serieschaltung}
\begin{parts}
  \part[1] Vergleichen Sie Ihre drei gemessenen Stromstärken
  untereinander. Was schliessen Sie daraus für den Strom in einer
  Serieschaltung?
  \begin{answerbox}
  \vspace{1.8cm}
  \end{answerbox}
  \begin{solution}
  Die drei Stromstärken sind (bis auf kleine Messabweichungen) gleich
  gross: In einer Serieschaltung ist der Strom überall gleich gross, da
  es nur einen einzigen Weg für die Ladungsträger gibt.
  \end{solution}

  \needspace{4.5cm}
  \part[2] Vergleichen Sie die Summe Ihrer drei gemessenen Spannungen
  mit $U_0$. Was schliessen Sie daraus?
  \begin{answerbox}
  \vspace{1.8cm}
  \end{answerbox}
  \begin{solution}
  Die Summe der drei Spannungen ($\SI{0.552}{\volt}+\SI{1.173}{\volt}
  +\SI{2.321}{\volt}=\SI{4.046}{\volt}$) stimmt mit $U_0=\SI{4}{\volt}$
  überein: Die angelegte Spannung wird auf die einzelnen Widerstände
  aufgeteilt. Als Formel:
  \[
    U_0 = U_1+U_2+U_3.
  \]
  \end{solution}

  \needspace{5cm}
  \part[2] Stellen Sie mit Ihrem gemessenen $R_{ges}$ und den drei
  Einzelwiderständen $R_1$, $R_2$, $R_3$ eine Formel für den
  Ersatzwiderstand $R_{ers}$ einer Serieschaltung auf: Ist er grösser
  oder kleiner als jeder einzelne Widerstand?
  \begin{answerbox}
  \vspace{2cm}
  \end{answerbox}
  \begin{solution}
  Der gemessene Gesamtwiderstand $R_{ges}=\SI{347}{\ohm}$ ist
  \textbf{\emph{grösser}} als jeder einzelne Widerstand: Jeder
  zusätzliche Widerstand bremst die Ladungsträger zusätzlich. Als Formel
  für den Ersatzwiderstand einer Serieschaltung:
  \[
    R_{ers} = R_1+R_2+R_3 = \SI{47}{\ohm}+\SI{100}{\ohm}+\SI{200}{\ohm}
    = \SI{347}{\ohm}.
  \]
  \end{solution}
\end{parts}

\needspace{10cm}
\question[6] \textbf{Parallelschaltung}
\begin{itemize}[leftmargin=1.2cm,itemsep=0.3em]
  \item Bauen Sie dieselben drei Widerstände gemäss Schaltplan parallel
  mit Schalter und Netzteil auf. Netzteil wieder auf $\SI{4}{\volt}$.
  \item Messen Sie die Spannung über jedem der drei Widerstände sowie
  die Gesamtspannung $U_0$.
  \item Messen Sie die Stromstärke $I$ durch jeden der drei Widerstände
  sowie die Gesamtstromstärke $I_0$, die aus dem Netzteil fliesst.
  \item Trennen Sie anschliessend die Spannungsquelle vom Stromkreis:
  Für eine Widerstandsmessung darf keine Spannungsquelle angeschlossen
  sein (siehe \glqq Widerstand messen\grqq{} weiter oben). Messen Sie
  dann mit dem Multimeter ($\Omega$-Modus) direkt den Gesamtwiderstand
  $R_{ges}$ derselben Parallelschaltung.
\end{itemize}
\begin{center}
\begin{circuitikz}[scale=1.15]
  \draw (0,0) to[battery1,invert,l=$U_0$] (0,3)
              to[normal open switch] (2,3);
  \draw (2,3) -- (8,3);
  \draw (0,0) -- (8,0);
  \draw (2,3) to[R,l=$R_1$] (2,0);
  \draw (5,3) to[R,l=$R_2$] (5,0);
  \draw (8,3) to[R,l=$R_3$] (8,0);
\end{circuitikz}
\end{center}
\begin{center}
\renewcommand{\arraystretch}{2.2}
\begin{tabular}{lcc}
\toprule
$R$ & $U$ in $\si{\volt}$ & $I$ in $\si{\ampere}$ \\
\midrule
$R_1$ & \rule{3cm}{0.4pt} & \rule{3cm}{0.4pt} \\
$R_2$ & \rule{3cm}{0.4pt} & \rule{3cm}{0.4pt} \\
$R_3$ & \rule{3cm}{0.4pt} & \rule{3cm}{0.4pt} \\
\bottomrule
\end{tabular}
\renewcommand{\arraystretch}{1}

\vspace{0.5cm}
Gesamtspannung $U_0 = $ \rule{3cm}{0.4pt} \hspace{1cm}
Gesamtstromstärke $I_0 = $ \rule{3cm}{0.4pt}

\vspace{0.3cm}
Gesamtwiderstand $R_{ges} = $ \rule{3cm}{0.4pt}
\end{center}
\begin{solution}
\Bild{7cm}{Bilder/Parallelschaltung.jpg}
\begin{center}
\renewcommand{\arraystretch}{2.2}
\begin{tabular}{lcc}
\toprule
$R$ & $U$ in $\si{\volt}$ & $I$ in $\si{\ampere}$ \\
\midrule
$R_1$ & $4$ & $0.086$ \\
$R_2$ & $4$ & $0.041$ \\
$R_3$ & $4$ & $0.02$ \\
\bottomrule
\end{tabular}
\renewcommand{\arraystretch}{1}

\vspace{0.5cm}
Gesamtspannung $U_0 = \SI{4}{\volt}$ \hspace{1cm}
Gesamtstromstärke $I_0 = \SI{0.148}{\ampere}$

\vspace{0.3cm}
Gesamtwiderstand $R_{ges} \approx \SI{27.6}{\ohm}$
\end{center}
\end{solution}

\needspace{6cm}
\question[5] \textbf{Konklusion: Parallelschaltung}
\begin{parts}
  \part[1] Vergleichen Sie Ihre drei gemessenen Spannungen
  untereinander. Was schliessen Sie daraus für die Spannung in einer
  Parallelschaltung?
  \begin{answerbox}
  \vspace{1.8cm}
  \end{answerbox}
  \begin{solution}
  Die drei Spannungen sind (bis auf kleine Messabweichungen) gleich
  gross: In einer Parallelschaltung liegt an allen Zweigen dieselbe
  Spannung an, da sie zwischen denselben zwei Punkten liegen.
  \end{solution}

  \needspace{4.5cm}
  \part[2] Vergleichen Sie die Summe Ihrer drei gemessenen
  Stromstärken mit $I_0$. Was schliessen Sie daraus?
  \begin{answerbox}
  \vspace{1.8cm}
  \end{answerbox}
  \begin{solution}
  Die Summe der drei Stromstärken ($\SI{0.086}{\ampere}
  +\SI{0.041}{\ampere}+\SI{0.02}{\ampere} = \SI{0.147}{\ampere}$) stimmt
  mit $I_0=\SI{0.148}{\ampere}$ überein: Der Gesamtstrom teilt sich auf
  die einzelnen Zweige auf. Als Formel:
  \[
    I_0 = I_1+I_2+I_3.
  \]
  \end{solution}

  \needspace{10cm}
  \part[2] Für den Ersatzwiderstand einer Parallelschaltung gilt die
  Formel
  \[
    \frac{1}{R_{ers}} = \frac{1}{R_1}+\frac{1}{R_2}+\frac{1}{R_3}.
  \]
  Überprüfen Sie mit Ihrem gemessenen $R_{ges}$ und den drei
  Einzelwiderständen, ob diese Formel stimmt. Ist $R_{ers}$ grösser oder
  kleiner als der kleinste der drei Einzelwiderstände?
  \begin{answerbox}
  \vspace{2cm}
  \end{answerbox}
  \begin{solution}
  Rechenweg: Zuerst die drei Widerstände in die Formel einsetzen:
  \[
    \frac{1}{R_{ers}} = \frac{1}{R_1}+\frac{1}{R_2}+\frac{1}{R_3}
    = \frac{1}{\SI{47}{\ohm}}+\frac{1}{\SI{100}{\ohm}}+\frac{1}{\SI{200}{\ohm}}.
  \]
  Jeden Bruch einzeln als Dezimalzahl ausrechnen:
  \[
    \frac{1}{R_{ers}} \approx \frac{0.02128}{\si{\ohm}}+\frac{0.01}{\si{\ohm}}
    +\frac{0.005}{\si{\ohm}}.
  \]
  Addieren:
  \[
    \frac{1}{R_{ers}} \approx \frac{0.03628}{\si{\ohm}}.
  \]
  Beide Seiten invertieren (Kehrwert bilden), um nach $R_{ers}$
  aufzulösen:
  \[
    R_{ers} \approx \frac{\si{\ohm}}{0.03628} \approx \SI{27.6}{\ohm}.
  \]
  Das stimmt mit dem gemessenen Gesamtwiderstand
  $R_{ges}\approx\SI{27.6}{\ohm}$ überein. $R_{ers}$ ist dabei
  \textbf{\emph{kleiner}} als der kleinste Einzelwiderstand
  ($R_1=\SI{47}{\ohm}$): Jeder zusätzliche Zweig bietet den
  Ladungsträgern einen weiteren Weg.
  \end{solution}
\end{parts}
\end{groupbox}
\end{questions}

% =============================================================================
% Merksaetze: Kirchhoffsche Regeln
% =============================================================================
\needspace{8cm}
\begin{theoriebox}[Merksätze]

\begin{formelbox}[Merksatz: Serieschaltung]
Bei einer Serieschaltung gibt es nur einen einzigen Weg für die
Ladungsträger: Der Strom ist überall gleich gross. Die von der
Spannungsquelle angelegte Spannung wird von den Bauteilen (den
\textbf{\emph{Wandlern}}) nacheinander umgewandelt: Die Summe der
Teilspannungen ergibt die Gesamtspannung.
\end{formelbox}

\begin{formelbox}[Merksatz: Parallelschaltung]
Bei einer Parallelschaltung liegen alle Bauteile zwischen denselben zwei
Punkten, also auf derselben Spannung: Wasser bzw. Strom auf gleicher
Höhe (gleicher Spannung) hat dasselbe \textbf{\emph{Wirkungspotential}},
deshalb liegt an allen parallelen Zweigen dieselbe Spannung $U_0$ an. An
jeder Verzweigung kommt genau so viel Strom hinzu, wie wegfliesst: Die
Summe der Teilströme ergibt den Gesamtstrom.
\end{formelbox}

\begin{formelbox}[Merksatz: Schaltungen allgemein]
Die beiden Regeln, die Sie oben für Serie- und Parallelschaltung
festgehalten haben, gelten so in jeder Schaltung, egal ob Serie,
Parallel oder eine Mischung aus beidem, und haben eigene Namen: Die
Stromregel an einer Verzweigung heisst \textbf{\emph{Knotenregel}}, die
Spannungsregel in der Schleife heisst \textbf{\emph{Maschenregel}}.
\end{formelbox}

\end{theoriebox}

% =============================================================================
% Ersatzwiderstand: Mischschaltung
% =============================================================================
\needspace{11cm}
\begin{questions}
\setcounter{question}{5}
\begin{groupbox}[Ersatzwiderstand: Mischschaltung]
\question[4] Die folgende Schaltung enthält Ihre drei Widerstände sowohl
in Serie als auch parallel geschaltet (wieder ohne Netzteil):
\begin{center}
\begin{circuitikz}[scale=1.7]
  \draw (0,1.5) to[R,l=$R_1$] (2.5,1.5);
  \filldraw (2.5,1.5) circle (0.05);
  \draw (2.5,1.5) -- (2.5,2.8) to[R,l=$R_2$] (5.5,2.8) -- (5.5,1.5);
  \draw (2.5,1.5) -- (2.5,0.2) to[R,l=$R_3$] (5.5,0.2) -- (5.5,1.5);
  \filldraw (5.5,1.5) circle (0.05);
  \draw (5.5,1.5) -- (7,1.5) -- (7,-1) -- (3.9,-1);
  \node[draw,circle,minimum size=0.8cm] (Om) at (3.5,-1) {$\Omega$};
  \draw (3.1,-1) -- (0,-1) -- (0,1.5);
\end{circuitikz}
\end{center}
Berechnen Sie den Ersatzwiderstand $R_{ers}$ dieser Schaltung mit den
Formeln von oben. Bauen Sie die Schaltung anschliessend auf dem
Steckbrett auf und messen Sie zur Kontrolle.
\begin{answerbox}
\vspace{2.5cm}
\end{answerbox}
\begin{solution}
Zuerst die Parallelschaltung von $R_2$ und $R_3$:
\[
  \frac{1}{R_{23}} = \frac{1}{\SI{100}{\ohm}}+\frac{1}{\SI{200}{\ohm}}
\]
Gemeinsamer Nenner $\SI{200}{\ohm}$:
\[
  \frac{1}{R_{23}} = \frac{2}{\SI{200}{\ohm}}+\frac{1}{\SI{200}{\ohm}}
  = \frac{3}{\SI{200}{\ohm}}
\]
Beide Seiten invertieren (Kehrwert bilden), um nach $R_{23}$ aufzulösen:
\[
  \Longrightarrow\quad R_{23} = \frac{\SI{200}{\ohm}}{3} \approx \SI{66.7}{\ohm}.
\]
Diese liegt in Serie mit $R_1$:
\[
  R_{ers} = R_1+R_{23} = \SI{47}{\ohm}+\SI{66.7}{\ohm} = \SI{113.7}{\ohm}.
\]
Der gemessene Wert sollte (bis auf kleine Abweichungen) damit
übereinstimmen.
\end{solution}
\end{groupbox}
\end{questions}

% =============================================================================
% Schaltungs-Tausch
% =============================================================================
\needspace{16cm}
\begin{questions}
\setcounter{question}{6}
\begin{groupbox}[Schaltungs-Tausch: Komplexere Schaltung berechnen]
\question[6]
\begin{parts}
  \part[2] Zeichnen Sie in Ihrer Gruppe eine eigene, etwas komplexere
  Schaltung mit mindestens vier Widerständen, die sowohl Serie- als auch
  Parallelteile enthält. Wählen Sie selbst Widerstandswerte, beschriften
  Sie diese in Ihrer Skizze, und geben Sie die Skizze an eine andere
  Gruppe weiter.
  \Antwortraum{5}

  \part[4] Berechnen Sie den Gesamtwiderstand der Schaltung, die Sie von
  der anderen Gruppe erhalten haben. Vergleichen Sie anschliessend Ihr
  Ergebnis mit der Gruppe, die die Schaltung gezeichnet hat.
  \begin{answerbox}
  \vspace{3cm}
  \end{answerbox}
\end{parts}

\begin{solution}
\textbf{Beispiel, wie eine getauschte Schaltung aussehen könnte:}
\begin{center}
\begin{circuitikz}[scale=0.7]
  \draw (1,1.5) -- (1,3) to[R,l=$R_a$] (3.5,3) to[R,l=$R_b$] (6,3) -- (6,1.5);
  \draw (1,1.5) -- (1,0) to[R,l=$R_c$] (6,0) -- (6,1.5);
  \filldraw (1,1.5) circle (0.05);
  \filldraw (6,1.5) circle (0.05);
  \draw (6,1.5) to[R,l=$R_d$] (8.5,1.5);
  \draw (8.5,1.5) -- (10,1.5) -- (10,-1.5) -- (5.9,-1.5);
  \node[draw,circle,minimum size=0.8cm] (Om) at (5.5,-1.5) {$\Omega$};
  \draw (5.1,-1.5) -- (0,-1.5) -- (0,1.5) -- (1,1.5);
\end{circuitikz}
\end{center}
Mit $R_a=\SI{100}{\ohm}$, $R_b=\SI{200}{\ohm}$, $R_c=\SI{300}{\ohm}$ und
$R_d=\SI{50}{\ohm}$: $R_a$ und $R_b$ liegen in Serie, diese Serie liegt
parallel zu $R_c$, und das Ganze liegt in Serie mit $R_d$.

Zuerst $R_a$ und $R_b$ (in Serie):
\[
  R_{ab} = R_a+R_b = \SI{100}{\ohm}+\SI{200}{\ohm} = \SI{300}{\ohm}.
\]
Dann $R_{ab}$ parallel zu $R_c$:
\[
  \frac{1}{R_{abc}} = \frac{1}{R_{ab}}+\frac{1}{R_c}
  = \frac{1}{\SI{300}{\ohm}}+\frac{1}{\SI{300}{\ohm}}
  = \frac{2}{\SI{300}{\ohm}}
  \quad\Longrightarrow\quad R_{abc} = \SI{150}{\ohm}.
\]
Zuletzt $R_{abc}$ in Serie mit $R_d$:
\[
  R_{ges} = R_{abc}+R_d = \SI{150}{\ohm}+\SI{50}{\ohm} = \SI{200}{\ohm}.
\]
\end{solution}
\end{groupbox}
\end{questions}

% =============================================================================
% Zusatzaufgabe: UND- und ODER-Schaltung
% =============================================================================
\needspace{10cm}
\begin{questions}
\setcounter{question}{7}
\begin{groupbox}[Zusatzaufgabe (für schnelle Gruppen): UND- und ODER-Schaltung]
Mit zwei Schaltern und einer Glühbirne lassen sich zwei ganz
unterschiedliche Verhalten bauen. Einmal müssen beide Schalter
geschlossen sein, damit die Glühbirne leuchtet: Die Schalter hängen
voneinander ab. Einmal genügt ein einzelner Schalter: Die Schalter
wirken unabhängig voneinander.

\needspace{10cm}
\question[4] \textbf{UND-Schaltung}
Bauen Sie mit dem Netzteil, zwei Schaltern und einer Glühbirne eine
Schaltung, bei der die Glühbirne nur leuchtet, wenn \textbf{\emph{beide}}
Schalter gleichzeitig geschlossen sind. Zeichnen Sie anschliessend selbst
den Schaltplan Ihrer Schaltung.
\Antwortraum{4}
\begin{solution}
\Bild{6cm}{Bilder/schaltung_und.jpg}
\begin{center}
\begin{circuitikz}[scale=0.9]
  \draw (0,0) to[battery1,invert,l=$U_0$] (0,3)
              to[normal open switch,l=$S_1$] (2.2,3)
              to[normal open switch,l=$S_2$] (4.4,3)
              to[lamp] (4.4,0)
              -- (0,0);
  \node at (0.4,1.85) {$+$};
  \node at (0.4,1.15) {$-$};
\end{circuitikz}
\end{center}
Die beiden Schalter liegen in \textbf{\emph{Serie}}: Erst wenn $S_1$
\textbf{\emph{und}} $S_2$ geschlossen sind, ist die Schleife
geschlossen und die Glühbirne leuchtet. Daher der Name
\glqq UND-Schaltung\grqq{}.
\end{solution}

\needspace{10cm}
\question[4] \textbf{ODER-Schaltung}
Bauen Sie eine Schaltung, bei der die Glühbirne bereits leuchtet, wenn
nur \textbf{\emph{ein einzelner}} Schalter geschlossen ist, unabhängig
davon, wie der andere Schalter steht. Zeichnen Sie auch hier selbst den
Schaltplan Ihrer Schaltung.
\Antwortraum{4}
\begin{solution}
\Bild{6cm}{Bilder/schaltung_oder.jpg}
\begin{center}
\begin{circuitikz}[scale=0.9]
  \draw (0,0) to[battery1,invert,l=$U_0$] (0,3) -- (1.5,3);
  \draw (1.5,3) to[normal open switch,l=$S_1$] (1.5,1.5);
  \draw (1.5,3) -- (3.5,3) to[normal open switch,l=$S_2$] (3.5,1.5);
  \draw (1.5,1.5) -- (2.5,1.5) -- (3.5,1.5);
  \draw (2.5,1.5) to[lamp] (2.5,0);
  \draw (2.5,0) -- (0,0);
  \node at (0.4,1.85) {$+$};
  \node at (0.4,1.15) {$-$};
\end{circuitikz}
\end{center}
Die beiden Schalter liegen \textbf{\emph{parallel}}: Schon $S_1$
\textbf{\emph{oder}} $S_2$ allein schliesst einen Weg für den Strom.
Daher der Name \glqq ODER-Schaltung\grqq{}.
\end{solution}
\end{groupbox}
\end{questions}

\end{document}
